% Pursuing Stacks (cotes 134-2 to 134-8): archive cote 134-6, site page 42 (PDF page 43), typescript page 360.
% typescript. Draft compared: ps5.tex:3437-3497.
% AI-transcribed by Claude Opus 5.5 (claude-opus-5-5, adaptive thinking, effort high), Message Batches API, batch ps3 (msgbatch_01XdFmuVVBAcgc7EFxdtCqsj), 2026-11-12; not yet checked by eye.
\scanpage{134-6/42}
\mspage{360}

then (X and) $X\setminus U = Y$ are aspheric too.

Proof of lemma 2: by an immediate induction, we are reduced to
the case when I ha\unsure{s} just two elements, $I = \hand{\{}1,\hand{2\}}$ , but then $X\setminus U$
can be viewed as the union of the two closed subcategories $X_1 x Y_2$ and
$X_2 x Y_1$ , whose intersection is $Y_1 x Y_2$ . As all three categories are
aspheric (being products of aspheric categories), it follows by the
well-known Mayer-Vietoris argument that so is $X\setminus U$ , qed.

Thus we are left with proving that $\hand{\Delta}^{f}$ is aspheric. Somewhat sur-%
prisingly, that's where I spent a number of hours not getting anywhere
and feeling foolish ! There is a very simple heuristic argument though
involving the standard calculation of the cohomology invariants of any
semisimplicial ``complex'', in terms of the standard boundary operation\hand{s:}
\unsure{\xout{+}} if we admit that the same calculations are valid when working with
semisimplicial ``face complexes'' i.e. objects of \xout{\illegible{1}} $\hand{\Delta}^{f\wedge}$, then it is
enough to apply this to the final object of $\hand{\Delta}^{f\wedge}$ (including for compu-%
\unsure{tation} \inserted{of} the non commutative $\unsure{\hand{H}}^1$ with constant coefficients) to get aspheri-%
city of $\hand{\Delta}^{f}$ . As a matter of fact, this argument would give directly
asphericity of $\hand{\alpha}$, bypassing altogether the categories $(*\ *)$ and lemmas
1 and 2 . Apparently I got a block against the down-to-earth computa-%
tional approach to cohomology via semisimplicial calculations, and
have been trying to bypass it at all price - and not succe\unsure{\hand{e}}ding !
Then curiously enough, when finding no other way out than look at tho-%
se boundary operations and try to understand what they meant (something
I remember vaguely have been doing \xout{already} \inserted[typed]{once} ages ago !), this brought
me back again to the abelianization story of sections 92 and 93, and
to a more comprehensive way for looking at \hand{``}abelianizators'' , and get
an existence and unicity statement for these. (At any rate, for a sui-%
tabl\over{e}{y} strengthened versi\xout{o}on of these.) This\xout{,} seems to me of independant
interest, and worth being written down with some care.

99. When writing down (in sections 92 and 93) some rambling reflexi\xout{ons}
ons about ``abelianisation'' and ``abelianizators'', there has been a
persistent feeling of uneasiness, which I kept pushing aside, as I
did'nt want to spend too much thought on this ``digression''. This unea-%
siness \xout{wasxtied} had \inserted{surely} something to do with the way abelianization \hand{(}of an
object X say of an elementary modelizer $A^{\wedge}$\hand{)} was handled\hand{,} \inserted{so that} it was desig-%
ned in a more or less exclusive way for embodying information about
the ``homology structure'' of the homotopy type modelized by X, or

\margin{42}
